% =====================================================================
%  Chapter: The relativistic point particle in electromagnetic and
%           non-abelian external fields
%  Conventions: Gaussian units, metric signature (+,-,-,-)
%
%  This file compiles stand-alone (pdflatex).  To embed it in a book,
%  copy the macros from the preamble and everything from \chapter
%  to the end of the "Notes and further reading" section.
% =====================================================================
\documentclass[11pt,a4paper]{book}

% (OT1 default fonts: scalable without extra packages)
\usepackage[utf8]{inputenc}
\usepackage{amsmath,amssymb,amsthm,mathtools,bm}
\usepackage{booktabs,array,enumitem}
\usepackage[margin=2.6cm]{geometry}
\usepackage[expansion=false]{microtype}
\usepackage{xcolor}
\usepackage{float}
\usepackage[colorlinks=true,linkcolor=blue!45!black,citecolor=blue!45!black,urlcolor=blue!45!black]{hyperref}

\numberwithin{equation}{section}

% ---------------------------------------------------------------- macros
\newcommand{\dd}{\mathrm{d}}
\newcommand{\ii}{\mathrm{i}}
\newcommand{\pdv}[2]{\dfrac{\partial #1}{\partial #2}}
\newcommand{\PB}[2]{\left\{#1,\,#2\right\}}
\newcommand{\tr}{\operatorname{tr}}
\newcommand{\Cc}{\mathcal{C}}          % mass-shell constraint
\newcommand{\bx}{\mathbf{x}}
\newcommand{\bv}{\mathbf{v}}
\newcommand{\bp}{\mathbf{p}}
\newcommand{\bA}{\mathbf{A}}
\newcommand{\bE}{\mathbf{E}}
\newcommand{\bB}{\mathbf{B}}
\newcommand{\bpi}{\boldsymbol{\pi}}

\theoremstyle{definition}
\newtheorem{problem}{Problem}[chapter]
\theoremstyle{remark}
\newtheorem*{remark}{Remark}

\begin{document}

% =====================================================================
\chapter[Relativistic Particle in Gauge Fields]{The Relativistic Point Particle in Electromagnetic and\\ Non-Abelian External Fields}
\label{ch:particle}
% =====================================================================

A relativistic point particle moving in a prescribed gauge field is the
simplest system in which Lorentz invariance, gauge invariance and the
canonical formalism meet.  Yet there is no single ``the'' Lagrangian for
it.  The same world lines follow from half a dozen different actions---some
regular but not manifestly covariant, some covariant but singular, some
covariant with an auxiliary field, some with a fifth ``evolution''
parameter---and each leads to its own Hamiltonian, its own constraints and
its own list of conserved quantities.  In this chapter we put these
descriptions side by side, first for the electromagnetic field and then for
a non-abelian field.  In the second case the particle has to carry an
internal ``colour'' degree of freedom, and the Lorentz force is replaced by
the equations of Wong.

Throughout, the field is \emph{external}: it is prescribed, the particle is
a test body, and both radiation reaction and the back-reaction of the
particle on the field are ignored.

\medskip
\noindent\textbf{Plan.}
Section~\ref{sec:conv} fixes conventions.  Section~\ref{sec:free} treats the
free particle, where the structure of the problem (reparametrization
invariance, the mass-shell constraint, the einbein) is easiest to see.
Section~\ref{sec:em} couples the particle to the electromagnetic field and
derives the Lorentz force from each formulation; it also introduces the
Stueckelberg--Horwitz--Piron (SHP) form with a fifth potential.
Section~\ref{sec:examples} solves three standard problems exactly.
Section~\ref{sec:nonabelian} generalizes to a non-abelian gauge field:
colour variables, Wong's equations, Hamiltonian formalism with the
Lie--Poisson bracket, the SHP generalization, and the Kaluza--Klein
geodesic picture.  Section~\ref{sec:summary} compares the formulations.
Problems and notes close the chapter.

% =====================================================================
\section{Conventions}\label{sec:conv}
% =====================================================================

\begin{itemize}[leftmargin=*]
\item \textbf{Units.}  Gaussian units; $c$ and $\hbar$ are kept explicit.
  The charge $q$ is measured in statcoulombs, and $q^2/\hbar c$ is
  dimensionless.

\item \textbf{Metric.}  $\eta_{\mu\nu}=\mathrm{diag}(+1,-1,-1,-1)$,
  $x^\mu=(ct,\bx)$, $a\!\cdot\! b=a^\mu b_\mu=a^0b^0-\mathbf a\!\cdot\!\mathbf b$,
  $\partial_\mu=\partial/\partial x^\mu=(c^{-1}\partial_t,\nabla)$.
  Greek indices run over $0,\dots,3$, Latin ones over $1,2,3$.
  A world line is $x^\mu(\lambda)$ with an \emph{arbitrary} parameter
  $\lambda$ and $\dot{}\equiv \dd/\dd\lambda$; the proper time is
  $\dd\tau=\dd s/c$, $\dd s^2=\eta_{\mu\nu}\dd x^\mu\dd x^\nu$.
  The four-velocity $u^\mu=\dd x^\mu/\dd\tau=(\gamma c,\gamma\bv)$ obeys
  $u\!\cdot\! u=c^2$, with $\gamma=(1-v^2/c^2)^{-1/2}$.

\item \textbf{Electromagnetic field.}
  $A^\mu=(\phi,\bA)$, $A_\mu=(\phi,-\bA)$,
  $\bE=-\nabla\phi-c^{-1}\partial_t\bA$, $\bB=\nabla\times\bA$, and
  \begin{equation}
    F_{\mu\nu}=\partial_\mu A_\nu-\partial_\nu A_\mu,\qquad
    F_{0i}=E^i,\quad F_{ij}=-\varepsilon_{ijk}B^k,\qquad
    F^{\mu\nu}=\begin{pmatrix}
      0&-E_x&-E_y&-E_z\\ E_x&0&-B_z&B_y\\ E_y&B_z&0&-B_x\\ E_z&-B_y&B_x&0
    \end{pmatrix}.
    \label{eq:F-components}
  \end{equation}

\item \textbf{Momenta and brackets.}
  For a Lagrangian $L(x,\dot x)$ the momentum conjugate to $x^\mu$ is
  $\partial L/\partial\dot x^\mu$.  With signature $(+,-,-,-)$ this has the
  ``wrong'' sign for the four-vector $p^\mu=(H/c,\bp)$ because
  $p_i=-p^i$.  Following Landau and Lifshitz we therefore \emph{define}
  \begin{equation}
    p_\mu\equiv-\pdv{L}{\dot x^\mu},\qquad p^\mu=(H/c,\bp),
    \label{eq:p-def}
  \end{equation}
  with $\bp$ the ordinary canonical momentum.  If $H_\lambda$ generates
  evolution in $\lambda$, Hamilton's equations read
  \begin{equation}
    \dot x^\mu=-\pdv{H_\lambda}{p_\mu},\qquad
    \dot p_\mu=\pdv{H_\lambda}{x^\mu},
    \label{eq:Hamilton-cov}
  \end{equation}
  and the Poisson bracket is
  $\PB fg=\partial_{p_\mu}f\,\partial_{x^\mu}g-\partial_{x^\mu}f\,\partial_{p_\mu}g$, so that
  \begin{equation}
    \PB{x^\mu}{p_\nu}=-\delta^\mu{}_\nu,\qquad
    \PB{x^i}{p^j}=\delta^{ij},\qquad
    \PB{p_\mu}{f(x)}=\partial_\mu f,\qquad \dot f=\PB f{H_\lambda}.
    \label{eq:PB-conv}
  \end{equation}
  The sign in $\PB{x^\mu}{p_\nu}$ is the price of the mostly-minus metric;
  the spatial part is the usual one.

\item \textbf{Kinetic momentum.}  $\pi_\mu\equiv p_\mu-\tfrac qc A_\mu$
  (and $\bpi=\bp-\tfrac qc\bA$).
\end{itemize}

% =====================================================================
\section{The free particle: one system, several actions}\label{sec:free}
% =====================================================================

\subsection{The geometric action and its coordinate-time form}

The world line of a free particle extremizes its proper length,
\begin{equation}
  S_0=-mc\int\!\dd s=-mc\int_{\lambda_1}^{\lambda_2}\!\sqrt{\dot x^\mu\dot x_\mu}\;\dd\lambda .
  \label{eq:S0}
\end{equation}
The integrand is homogeneous of degree one in $\dot x^\mu$, hence
$S_0$ is invariant under any reparametrization $\lambda\to\lambda'(\lambda)$.
Choosing $\lambda=t$ gives the familiar \emph{non-covariant} Lagrangian
\begin{equation}
  L_0=-mc^2\sqrt{1-v^2/c^2}=-mc^2+\tfrac12mv^2+\tfrac18\,\frac{mv^4}{c^2}+\dots
  \label{eq:L0-t}
\end{equation}
with
\begin{equation}
  \bp=\pdv{L_0}{\bv}=\gamma m\bv,\qquad
  H_0=\bp\!\cdot\!\bv-L_0=\gamma mc^2=c\sqrt{\bp^2+m^2c^2}.
  \label{eq:H0-t}
\end{equation}
This Lagrangian is \emph{regular}: the Hessian
$\partial^2L_0/\partial v^i\partial v^j=\gamma m(\delta_{ij}+\gamma^2v_iv_j/c^2)$
is positive definite, so the Legendre transform is invertible and no
constraints arise.  Regularity is bought at the price of manifest covariance:
time has been singled out.

\subsection{Manifest covariance and the mass-shell constraint}

Keep $\lambda$ arbitrary.  The covariant momentum is
\begin{equation}
  p_\mu=-\pdv{L}{\dot x^\mu}=\frac{mc\,\dot x_\mu}{\sqrt{\dot x^2}}
  \quad\Longrightarrow\quad
  p_\mu p^\mu=m^2c^2 \ \ \text{identically}.
  \label{eq:massshell-free}
\end{equation}
The four components of $p_\mu$ are not independent functions of the
velocities: the relativistic dispersion relation $E^2=p^2c^2+m^2c^4$ is a
\emph{primary constraint}.  The cause is homogeneity.  The Hessian
\begin{equation}
  \pdv{^2L}{\dot x^\mu\partial\dot x^\nu}
  =-\frac{mc}{\sqrt{\dot x^2}}\Bigl(\eta_{\mu\nu}-\frac{\dot x_\mu\dot x_\nu}{\dot x^2}\Bigr)
\end{equation}
annihilates $\dot x^\nu$, so it is singular, and Euler's theorem
$\dot x^\mu\,\partial L/\partial\dot x^\mu=L$ shows that the canonical
Hamiltonian vanishes identically,
\begin{equation}
  H_{\rm can}=-p_\mu\dot x^\mu-L=0 .
\end{equation}
In Dirac's language the dynamics is generated entirely by constraints:
\begin{equation}
  H_T=e(\lambda)\,\Cc,\qquad
  \Cc\equiv\tfrac12\bigl(m^2c^2-p_\mu p^\mu\bigr)\approx0,
  \label{eq:HT-free}
\end{equation}
with an arbitrary multiplier $e(\lambda)$.  Hamilton's equations
\eqref{eq:Hamilton-cov} give $\dot x^\mu=e\,p^\mu$ and $\dot p_\mu=0$:
straight lines with arbitrary parametrization.  The constraint is
first class (there is only one, so $\PB{\Cc}{\Cc}=0$), and its gauge
transformations \emph{are} the reparametrizations.

Fixing the gauge breaks the symmetry and recovers the regular formulation.
Choose $\chi\equiv x^0-c\lambda\approx0$.  Since
$\PB{\chi}{\Cc}=p^0\neq0$ on shell, the pair $(\Cc,\chi)$ is second class;
solving $\Cc=0$ for $p^0$ gives $p^0=\pm\sqrt{\bp^2+m^2c^2}$, and the
generator of evolution in $t=\lambda$ is $H=c\,p^0=\pm c\sqrt{\bp^2+m^2c^2}$.
The positive root is \eqref{eq:H0-t}; the negative root describes world
lines traversed backwards in time, which in the Feynman--Stueckelberg
interpretation are antiparticles.

\subsection{The einbein (Polyakov) form}

The square root in \eqref{eq:S0} is awkward and does not allow $m=0$.  Introduce an
auxiliary field $e(\lambda)$, the \emph{einbein}:
\begin{equation}
  S_e[x,e]=-\frac12\int\!\dd\lambda\Bigl(\frac{\dot x^\mu\dot x_\mu}{e}+e\,m^2c^2\Bigr).
  \label{eq:Se-free}
\end{equation}
The equation of motion for $e$ is algebraic, $\dot x^2=e^2m^2c^2$, so
$e=\sqrt{\dot x^2}/(mc)$, and substituting back returns \eqref{eq:S0}.
Under $\lambda\to\lambda'$ the einbein transforms as a density,
$e'\,\dd\lambda'=e\,\dd\lambda$, so $e\,\dd\lambda=\dd\tau/m$ is the invariant
proper time in units of $1/m$.  Now $p_\mu=\dot x_\mu/e$ and
\begin{equation}
  H=-p_\mu\dot x^\mu-L=-\frac{\dot x^2}{e}+\frac12\Bigl(\frac{\dot x^2}{e}+em^2c^2\Bigr)
  =e\,\Cc,
\end{equation}
which is exactly the total Hamiltonian \eqref{eq:HT-free}: the einbein is
the Lagrange multiplier of the mass-shell constraint.  (In Dirac's
analysis $e$ has a vanishing conjugate momentum $p_e\approx0$ as a primary
constraint, and its preservation produces $\Cc\approx0$ as a secondary
constraint.)  The gauge $e=1/m$, $\lambda=\tau$ gives
$L=-\tfrac12m(\dot x^2+c^2)$, and the constraint becomes the normalization
$\dot x^2=c^2$.  For $m=0$, $S_e=-\tfrac12\int\dot x^2/e\,\dd\lambda$ describes a
light ray.

\subsection{The quadratic (Stueckelberg) Lagrangian}

If one freezes the einbein, $e=1/M$ with $M$ a fixed constant, and drops the
constant term, one obtains
\begin{equation}
  L_{S}=-\frac M2\,\dot x^\mu\dot x_\mu,\qquad
  p_\mu=M\dot x_\mu,\qquad
  K=-p_\mu\dot x^\mu-L_S=-\frac{1}{2M}\,p_\mu p^\mu .
  \label{eq:LS-free}
\end{equation}
This is regular, quadratic in the velocities, and involves no constraint.
It is the Lagrangian of Stueckelberg's parametrized relativistic mechanics
(Section~\ref{sec:shp}).  The price is that the evolution parameter $\tau$
is \emph{not} the proper time, and $\dot x^2$ is not fixed by the
equations of motion but by initial conditions: for a free particle $K$ is
conserved, and fixing $K=-m^2c^2/2M$ reproduces the mass shell.  Mass is
promoted from a parameter to a constant of motion.  (In the $(-,+,+,+)$
literature $K=+p^2/2M$ appears; the sign difference is the signature.)

% =====================================================================
\section{The charged particle in an electromagnetic field}\label{sec:em}
% =====================================================================

\subsection{Minimal coupling}

The simplest Lorentz-scalar, reparametrization-invariant, local coupling of
a world line to a one-form $A_\mu$ is $\int A_\mu\dd x^\mu$:
\begin{equation}
  S=-mc\int\!\dd s-\frac qc\int\!A_\mu\,\dd x^\mu
   =\int\!\dd\lambda\Bigl(-mc\sqrt{\dot x^2}-\frac qc\,A_\mu\dot x^\mu\Bigr).
  \label{eq:S-EM}
\end{equation}
With $\lambda=t$, $A_\mu\dot x^\mu=c\phi-\bA\!\cdot\!\bv$, and
\begin{equation}
  L=-mc^2\sqrt{1-v^2/c^2}-q\phi+\frac qc\,\bv\!\cdot\!\bA .
  \label{eq:L-EM-t}
\end{equation}
For $q>0$ and static $\phi$ the potential energy is $q\phi$, as it should be.

\subsection{Canonical momentum and Hamiltonian (coordinate time)}

From \eqref{eq:L-EM-t},
\begin{equation}
  \bp=\pdv L\bv=\gamma m\bv+\frac qc\bA\equiv\bpi+\frac qc\bA ,
  \label{eq:p-EM}
\end{equation}
and
\begin{equation}
  H=\bp\!\cdot\!\bv-L=\gamma mc^2+q\phi
  =\boxed{\,c\sqrt{\bigl(\bp-\tfrac qc\bA\bigr)^2+m^2c^2}+q\phi\,}\ .
  \label{eq:H-EM}
\end{equation}
(The velocity-dependent terms cancel in $\bp\!\cdot\!\bv-L$ because the
magnetic coupling is linear in $\bv$.)  For $|\bpi|\ll mc$ this reduces to
$mc^2+\bpi^2/2m+q\phi$, the nonrelativistic Hamiltonian with
$\bp\to\bp-q\bA/c$.

Hamilton's equations give $\bv=\partial H/\partial\bp=c\bpi/\sqrt{\bpi^2+m^2c^2}$ and
\begin{equation}
  \dot p_i=-\pdv H{x^i}=-q\,\partial_i\phi+\frac qc\,v^j\partial_iA^j .
\end{equation}
Since $\bpi=\bp-q\bA/c$ and $\dd A^i/\dd t=\partial_tA^i+v^j\partial_jA^i$,
\begin{equation}
  \dot\pi^i=q\Bigl(-\partial_i\phi-\frac1c\partial_tA^i\Bigr)
  +\frac qc\,v^j\bigl(\partial_iA^j-\partial_jA^i\bigr)
  =q\bigl[\bE+\tfrac1c\,\bv\times\bB\bigr]^i ,
  \label{eq:Lorentz-3d}
\end{equation}
where we used $[\bv\times(\nabla\times\bA)]^i=v^j\partial_iA^j-v^j\partial_jA^i$.
This is the Lorentz force, with $\bpi=\gamma m\bv$.

\subsection{Covariant equations of motion}

With arbitrary $\lambda$, the Euler--Lagrange equations of
\eqref{eq:S-EM} read $\dd p_\mu/\dd\lambda=\tfrac qc\,\dot x^\nu\partial_\mu A_\nu$ with
\begin{equation}
  p_\mu=\frac{mc\,\dot x_\mu}{\sqrt{\dot x^2}}+\frac qc A_\mu .
  \label{eq:p-cov-EM}
\end{equation}
In the gauge $\lambda=\tau$ the kinetic part is $m u_\mu$, and
$\dd A_\mu/\dd\tau=u^\nu\partial_\nu A_\mu$ turns the equation into
\begin{equation}
  \boxed{\,m\,\frac{\dd u^\mu}{\dd\tau}=\frac qc\,F^{\mu\nu}u_\nu\,}
  \qquad\Longleftrightarrow\qquad
  \frac{\dd\bpi}{\dd t}=q\Bigl(\bE+\frac{\bv}c\times\bB\Bigr),\quad
  \frac{\dd(\gamma mc^2)}{\dd t}=q\,\bE\!\cdot\!\bv .
  \label{eq:Lorentz-cov}
\end{equation}
Antisymmetry of $F^{\mu\nu}$ gives $u_\mu\,\dd u^\mu/\dd\tau=0$, consistent with
$u^2=c^2$: the field never changes the particle's mass.  The magnetic
field does no work.  (The component form follows from
\eqref{eq:F-components}: e.g.\ $m\,\dd u^i/\dd\tau=\tfrac qc(F^{i0}u_0+F^{ij}u_j)
=q\gamma(E^i+(\bv\times\bB)^i/c)$.)

\subsection{Gauge invariance}

Under $A_\mu\to A_\mu+\partial_\mu\Lambda$, that is,
$\phi\to\phi+c^{-1}\partial_t\Lambda$, $\bA\to\bA-\nabla\Lambda$,
\begin{equation}
  L\to L-\frac qc\,\frac{\dd\Lambda}{\dd t},\qquad
  S\to S-\frac qc\bigl[\Lambda\bigr]_{\rm endpoints}.
\end{equation}
The Euler--Lagrange equations are unchanged.  The canonical variables are
\emph{not} invariant: $\bp\to\bp-\tfrac qc\nabla\Lambda$ and
$H\to H+\tfrac qc\partial_t\Lambda$, while $\bpi$, the kinetic energy
$\gamma mc^2$, and the force are invariant.  The transformation of
$(\bx,\bp,H)$ is a (time-dependent) canonical transformation.  Physical
observables are gauge-invariant functions of $(\bx,\bpi)$.  In the quantum
theory the same transformation is a phase,
$\psi\to e^{-\ii q\Lambda/\hbar c}\psi$.

\subsection{Symmetries and conserved quantities}

Let $\xi^\mu(x)$ be a Killing vector of Minkowski space (a translation,
rotation or boost) such that the field is invariant up to a gauge
transformation, $\mathcal L_\xi A_\mu=\partial_\mu\Lambda_\xi$.  Then
$\delta x^\mu=\varepsilon\xi^\mu$ changes $L$ by a total derivative
$-\tfrac qc\varepsilon\,\dd\Lambda_\xi/\dd\lambda$, and Noether's theorem gives the
conserved quantity
\begin{equation}
  \xi^\mu p_\mu-\frac qc\,\Lambda_\xi=\text{const}.
  \label{eq:Noether-EM}
\end{equation}
Examples: for a translation $n^\mu$ with $n^\mu\partial_\mu A_\nu=0$ the
\emph{canonical} (not kinetic) momentum component $n^\mu p_\mu$ is
conserved; for rotation about $z$ in a field with $\bA=\tfrac12\bB\times\bx$
the canonical angular momentum $p_\varphi=\gamma m\rho^2\dot\varphi+\tfrac{qB}{2c}\rho^2$ is
conserved.

\subsection{Covariant Hamiltonian formalism}\label{sec:cov-ham}

The constraint structure of the free particle carries over with
$p_\mu\to\pi_\mu=p_\mu-\tfrac qcA_\mu$:
\begin{equation}
  \pi_\mu\pi^\mu=m^2c^2,\qquad
  H_T=e\,\Cc,\quad
  \Cc=\tfrac12\bigl(m^2c^2-\pi_\mu\pi^\mu\bigr)\approx0 .
  \label{eq:HT-EM}
\end{equation}
(This is also the Hamiltonian of the einbein action
\begin{equation}
  S_e=\int\!\dd\lambda\Bigl[-\frac{\dot x^2}{2e}-\frac e2m^2c^2-\frac qc A_\mu\dot x^\mu\Bigr],
  \label{eq:Se-EM}
\end{equation}
for which $p_\mu=\dot x_\mu/e+\tfrac qc A_\mu$.)  Hamilton's equations
\eqref{eq:Hamilton-cov} give
\begin{equation}
  \dot x^\mu=e\,\pi^\mu,\qquad
  \dot p_\mu=\frac qc\,e\,\pi^\nu\partial_\mu A_\nu=\frac qc\,\dot x^\nu\partial_\mu A_\nu,
\end{equation}
which is the Euler--Lagrange equation above.  A particularly clean way to
see the Lorentz force is through the brackets of the \emph{kinetic}
momenta.  From \eqref{eq:PB-conv},
\begin{equation}
  \PB{x^\mu}{x^\nu}=0,\qquad
  \PB{x^\mu}{\pi_\nu}=-\delta^\mu{}_\nu,\qquad
  \boxed{\,\PB{\pi_\mu}{\pi_\nu}=-\frac qc\,F_{\mu\nu}\,}
  \label{eq:PB-pi-pi}
\end{equation}
(for the spatial components, $\PB{\pi^i}{\pi^j}=\tfrac qc\,\varepsilon_{ijk}B^k$).
Then
\begin{equation}
  \dot\pi_\mu=\PB{\pi_\mu}{e\Cc}=-e\,\pi^\nu\PB{\pi_\mu}{\pi_\nu}
  =\frac qc F_{\mu\nu}\,e\pi^\nu=\frac qc F_{\mu\nu}\dot x^\nu .
\end{equation}
The algebra \eqref{eq:PB-pi-pi}, together with $\dot x^\mu\propto\pi^\mu$,
\emph{is} the Lorentz force; the field strength appears as the curvature
(non-commutativity) of the kinetic momenta, a point of view that
generalizes verbatim to the non-abelian case and underlies Feynman's
derivation of the Lorentz force from commutation relations (Dyson, 1990).

\paragraph{Reduction to coordinate time.}
Solve $\Cc=0$ for $\pi^0=(p^0-q\phi/c)$ with the gauge $x^0=c\lambda$:
$(H-q\phi)^2/c^2=\bpi^2+m^2c^2$, i.e.\ \eqref{eq:H-EM}.

\paragraph{First-order (phase-space) action.}
Equivalently, vary independently $x^\mu$, $p_\mu$ and $e$ in
\begin{equation}
  S[x,p,e]=-\int\!\dd\lambda\,\bigl(p_\mu\dot x^\mu+e\,\Cc\bigr)
  \ \xrightarrow{\ e\Cc=0\ }\ 
  \int\bigl(\bp\!\cdot\!\dd\bx-H\,\dd t\bigr),
  \label{eq:S-first-order}
\end{equation}
where we used $-p_\mu\dot x^\mu=-p_0\dot x^0-p_i\dot x^i=\bp\!\cdot\!\dot\bx-H\dot t$.
Varying $p_\mu$ gives $\dot x^\mu=e\pi^\mu$, varying $e$ gives
$\Cc=0$; eliminating $p_\mu$ reproduces \eqref{eq:Se-EM}.

\subsection{The Hamilton--Jacobi equation}

Since $p_\mu=-\partial_\mu S$ for the on-shell action, the constraint
\eqref{eq:HT-EM} becomes the Hamilton--Jacobi equation
\begin{equation}
  \Bigl(\partial_\mu S+\frac qcA_\mu\Bigr)\Bigl(\partial^\mu S+\frac qcA^\mu\Bigr)=m^2c^2 .
  \label{eq:HJ-EM}
\end{equation}
It is the eikonal limit of the Klein--Gordon equation
$\bigl[D_\mu D^\mu+m^2c^2/\hbar^2\bigr]\psi=0$, $D_\mu=\partial_\mu+\ii\tfrac q{\hbar c}A_\mu$,
with $\psi=e^{\ii S/\hbar}$.  For the free particle $S=-p_\mu x^\mu$.

\subsection{Stueckelberg--Horwitz--Piron mechanics with a fifth potential}\label{sec:shp}

In the SHP framework (Stueckelberg 1941; Horwitz and Piron 1973) events are
parametrized by a universal evolution parameter $\tau$, with the dimension
of time and independent of the proper time; the mass of the particle is not
fixed a priori.  The dynamics uses the quadratic Lagrangian
\eqref{eq:LS-free}.  If the fields are allowed to depend on $\tau$, gauge
invariance under $A_\mu\to A_\mu+\partial_\mu\Lambda(x,\tau)$ requires a
fifth potential $A_5(x,\tau)$.  With $x^5\equiv c\tau$ (so
$\partial_5=c^{-1}\partial_\tau$, $\dot x^5=c$),
\begin{equation}
  L_{\rm SHP}=-\frac M2\,\dot x^\mu\dot x_\mu-\frac qc\bigl(A_\mu\dot x^\mu+c\,A_5\bigr),
  \qquad A_5\to A_5+\partial_5\Lambda ,
  \label{eq:L-SHP}
\end{equation}
because $A_\mu\dot x^\mu+cA_5\to A_\mu\dot x^\mu+cA_5+\dd\Lambda/\dd\tau$.  Then
\begin{equation}
  p_\mu=M\dot x_\mu+\frac qc A_\mu,\qquad
  K=-\frac1{2M}\,\pi_\mu\pi^\mu+q\,A_5,
  \label{eq:K-SHP}
\end{equation}
and the Euler--Lagrange equations are
\begin{equation}
  M\ddot x_\mu=\frac qc\,F_{\mu\nu}\dot x^\nu+q\,F_{\mu5},\qquad
  F_{\mu5}=\partial_\mu A_5-\frac1c\,\partial_\tau A_\mu .
  \label{eq:EOM-SHP}
\end{equation}
Two things are new compared with Maxwell dynamics.  First, the ``mass''
$m\equiv M\sqrt{\dot x^2}/c$ is no longer constant:
\begin{equation}
  \frac{\dd}{\dd\tau}\Bigl(\tfrac12M\dot x^2\Bigr)=q\,\dot x^\mu F_{\mu5},
  \qquad
  m^2c^2=M^2\dot x^2=2M\bigl(qA_5-K\bigr)\quad(\partial_\tau K=0).
  \label{eq:mass-SHP}
\end{equation}
Second, in the ``concatenation'' limit $A_5=0$, $\partial_\tau A_\mu=0$,
the extra force vanishes, $\dot x^2$ is constant, and
\eqref{eq:EOM-SHP} reduces to the Lorentz force \eqref{eq:Lorentz-cov} with
$\tau\propto$ proper time and $m=M\sqrt{\dot x^2}/c$.

% ---------------------------------------------------------------------
\subsection{The formulations at a glance}
% ---------------------------------------------------------------------

\begin{table}[H]
\centering\small
\renewcommand{\arraystretch}{1.5}
\begin{tabular}{@{}>{\raggedright\arraybackslash}p{2.7cm}>{\raggedright\arraybackslash}p{6.2cm}>{\raggedright\arraybackslash}p{5.7cm}@{}}
\toprule
Formulation & Lagrangian / action & Hamiltonian and constraints\\
\midrule
Coordinate time (regular)
 & $L=-\dfrac{mc^2}{\gamma}-q\phi+\dfrac qc\,\bv\!\cdot\!\bA$
 & $H=c\sqrt{\bpi^2+m^2c^2}+q\phi$; no constraints; not manifestly covariant.\\
Geometric (square root)
 & $L=-mc\sqrt{\dot x^2}-\dfrac qc A\!\cdot\!\dot x$
 & $H_{\rm can}=0$; $\Cc=\tfrac12(m^2c^2-\pi^2)\approx0$; singular, covariant.\\
Einbein
 & $L=-\dfrac{\dot x^2}{2e}-\dfrac{e\,m^2c^2}2-\dfrac qc A\!\cdot\!\dot x$
 & $H_T=e\Cc$; $e$ = multiplier; valid for $m=0$.\\
First-order
 & $S=-\!\int(p_\mu\dot x^\mu+e\Cc)\,\dd\lambda$
 & Independent $x^\mu,p_\mu,e$; reduces to $\int(\bp\!\cdot\!\dd\bx-H\,\dd t)$.\\
Quadratic (SHP)
 & $L=-\dfrac M2\dot x^2-\dfrac qc(A\!\cdot\!\dot x+cA_5)$
 & $K=-\dfrac{\pi^2}{2M}+qA_5$; no constraint; $m$ not fixed.\\
5D geodesic (\S\ref{sec:KK})
 & $L=-\dfrac M2\,G_{MN}\dot x^M\dot x^N$
 & Charge = momentum of a cyclic coordinate.\\
\bottomrule
\end{tabular}
\caption{Equivalent descriptions of a charged relativistic particle.  The
first two rows describe the same dynamics with different gauge choices
(time as parameter vs.\ arbitrary parameter); the first row is the gauge-fixed
version of the second.}
\label{tab:formulations}
\end{table}

% =====================================================================
\section{Three exact solutions}\label{sec:examples}
% =====================================================================

\subsection{Uniform electric field: hyperbolic motion}

Let $\bE=E\hat{\mathbf x}$ be constant and the particle start at rest at the
origin.  With $a\equiv qE/m$ the equations \eqref{eq:Lorentz-cov} for
$(u^0,u^1)$ are
\begin{equation}
  \frac{\dd u^0}{\dd\tau}=\frac ac\,u^1,\qquad
  \frac{\dd u^1}{\dd\tau}=\frac ac\,u^0,
\end{equation}
and, with $u^2=c^2$,
\begin{equation}
  u^0=c\cosh\frac{a\tau}c,\quad
  u^1=c\sinh\frac{a\tau}c,\qquad
  ct=\frac{c^2}a\sinh\frac{a\tau}c,\quad
  x=\frac{c^2}a\Bigl(\cosh\frac{a\tau}c-1\Bigr).
  \label{eq:hyperbolic}
\end{equation}
Eliminating $\tau$: $(x+c^2/a)^2-(ct)^2=(c^2/a)^2$, a hyperbola with proper
acceleration $a$, and
\begin{equation}
  x(t)=\frac{mc^2}{qE}\Bigl[\sqrt{1+\bigl(qEt/mc\bigr)^2}-1\Bigr]
  \ \longrightarrow\ \tfrac12at^2\ \ (t\ll mc/qE).
\end{equation}
In the Hamiltonian formalism with $\phi=-Ex$, $\bA=0$:
$H=c\sqrt{p_x^2+m^2c^2}-qEx$, so $\dot p_x=qE$ and $p_x=qEt$.

\subsection{Uniform magnetic field}

For $\bB=B\hat{\mathbf z}$ the speed, hence $\gamma$, is constant, and
$\gamma m\,\dot\bv=\tfrac qc\bv\times\bB$ gives uniform circular motion in the
plane perpendicular to $\bB$ with
\begin{equation}
  \omega_c=\frac{qB}{\gamma mc},\qquad
  R=\frac{v_\perp}{\omega_c}=\frac{p_\perp c}{qB},
\end{equation}
and uniform drift along $\bB$.  In the symmetric gauge the conserved
quantity is the canonical angular momentum of \eqref{eq:Noether-EM}.

\subsection{Plane wave}

Let $A^\mu=A^\mu(\varphi)$ with $\varphi=k_\mu x^\mu$, $k^2=0$ and
$k\!\cdot\!A=0$ (radiation gauge), so that
$F_{\mu\nu}=k_\mu A'_\nu-k_\nu A'_\mu$, with $'\equiv\dd/\dd\varphi$.  Because
$k^\mu F_{\mu\nu}=0$,
\begin{equation}
  \frac{\dd(k\!\cdot\!u)}{\dd\tau}=0,\qquad
  \frac{\dd\varphi}{\dd\tau}=k\!\cdot\!u_0\equiv\text{const},
\end{equation}
and the equation of motion becomes linear in $\varphi$:
$\dd u_\mu/\dd\varphi=-\tfrac q{mc}A'_\mu+\tfrac q{mc\,(k\cdot u_0)}\,k_\mu(A'\!\cdot\!u)$.
Its solution with $u^\mu(\varphi_0)=u_0^\mu$ and $a^\mu\equiv A^\mu(\varphi)-A^\mu(\varphi_0)$ is
\begin{equation}
  u^\mu(\varphi)=u_0^\mu-\frac q{mc}\,a^\mu
  +k^\mu\Bigl[\frac{q\,(u_0\!\cdot\!a)}{mc\,(k\!\cdot\!u_0)}-\frac{q^2\,a\!\cdot\!a}{2m^2c^2\,(k\!\cdot\!u_0)}\Bigr],
  \label{eq:Volkov-classical}
\end{equation}
as one verifies by substitution (the coefficient of $k^\mu$ is fixed by
$u^2=c^2$ together with $k\!\cdot\!a=0$).  Since $a\!\cdot\!a=-|\mathbf a|^2\le0$ in
radiation gauge, a particle initially at rest has
\begin{equation}
  \gamma(\varphi)=1+\frac{q^2}{2m^2c^4}\,\bigl|\mathbf A(\varphi)-\mathbf A(\varphi_0)\bigr|^2 ,
\end{equation}
so that the energy it carries away depends only on the instantaneous
vector potential relative to its initial value---a classical
``Volkov'' solution, with the $k^\mu$ term responsible for the drift along
the propagation direction.

% =====================================================================
\section{Non-abelian external fields}\label{sec:nonabelian}
% =====================================================================

\subsection{Gauge field and covariant derivative}

Let $G$ be a compact Lie group with hermitian generators $T^a$ in a
representation $R$ of dimension $d$,
\begin{equation}
  [T^a,T^b]=\ii f^{abc}T^c,\qquad f^{abc}\ \text{real, totally antisymmetric},
\end{equation}
(for example $T^a=\sigma^a/2$, $f^{abc}=\varepsilon^{abc}$ for $SU(2)$).
Colour indices $a,b,c$ are summed when repeated.  The gauge potential is
$\mathcal A_\mu=A^a_\mu T^a$, and
\begin{equation}
  D_\mu=\partial_\mu+\ii\lambda\,\mathcal A_\mu,\qquad
  \mathcal F_{\mu\nu}=\frac1{\ii\lambda}[D_\mu,D_\nu]
  =\partial_\mu\mathcal A_\nu-\partial_\nu\mathcal A_\mu+\ii\lambda[\mathcal A_\mu,\mathcal A_\nu],
  \label{eq:D-F}
\end{equation}
\begin{equation}
  F^a_{\mu\nu}=\partial_\mu A^a_\nu-\partial_\nu A^a_\mu-\lambda f^{abc}A^b_\mu A^c_\nu .
  \label{eq:Fa}
\end{equation}
The sign of $\lambda$ is chosen so that for a single abelian generator
$\lambda\mathcal A_\mu\to\tfrac q{\hbar c}A_\mu$, matching the electromagnetic
covariant derivative of Section~\ref{sec:em}.  In Gaussian units $A^a_\mu$
has the dimension of an electrostatic potential, and $\lambda\equiv g/\hbar c$,
where $g$ is the Yang--Mills coupling constant, with the dimension of an
electric charge ($g^2/\hbar c$ is dimensionless).  The classical equations
depend on $g$ and $\hbar$ only through the combination $\lambda$; the
$\hbar$ in $\lambda$ is a matter of normalization of the structure constants,
not a quantum effect.

Under a gauge transformation $U(x)\in G$,
\begin{equation}
  \psi\to U\psi,\qquad
  \mathcal A_\mu\to U\mathcal A_\mu U^{-1}+\frac\ii\lambda\,(\partial_\mu U)\,U^{-1},\qquad
  \mathcal F_{\mu\nu}\to U\mathcal F_{\mu\nu}U^{-1},
  \label{eq:gauge-NA}
\end{equation}
and for $U=1+\ii\alpha^aT^a$,
\begin{equation}
  \delta A^a_\mu=-\frac1\lambda\,(D_\mu\alpha)^a,\qquad
  (D_\mu\alpha)^a=\partial_\mu\alpha^a-\lambda f^{abc}A^b_\mu\alpha^c .
  \label{eq:gauge-NA-inf}
\end{equation}
The field strength is no longer invariant but rotates in the adjoint
representation.

\subsection{Why the charge must become dynamical}

In electrodynamics $q$ is a constant because $U(1)$ is abelian.  If a
particle were to couple to $F^a_{\mu\nu}$ through a \emph{fixed} colour
vector $Q^a$, the force $Q^aF^a_{\mu\nu}u^\nu$ would not be gauge invariant,
since $F^a$ rotates by a position-dependent adjoint rotation along the
world line.  Gauge invariance of the physical force therefore requires $Q^a$
to rotate together with $F^a$, and the rotation depends on the path.  So
$Q^a(\tau)$ must itself be a dynamical variable obeying a
\emph{covariant constancy} equation, $DQ^a/\dd\tau=0$.  The simplest
Lagrangian that realizes this is obtained by giving the particle an
internal state $\psi(\tau)\in\mathbb C^d$ in $R$.

\subsection{Colour variables and Wong's equations}\label{sec:wong}

Consider
\begin{equation}
  L=-mc\sqrt{\dot x^2}+\ii\hbar\,\psi^\dagger D_\lambda\psi,\qquad
  D_\lambda\psi\equiv\dot\psi+\ii\lambda\,\dot x^\mu\mathcal A_\mu\psi,\qquad
  \psi^\dagger\psi=1 .
  \label{eq:L-wong}
\end{equation}
Up to a total derivative (which makes $L$ real) this is
\begin{equation}
  L=-mc\sqrt{\dot x^2}+\frac{\ii\hbar}2\bigl(\psi^\dagger\dot\psi-\dot\psi^\dagger\psi\bigr)
  -\frac1c\,Q^aA^a_\mu\dot x^\mu,\qquad
  Q^a\equiv g\,\psi^\dagger T^a\psi,
  \label{eq:L-wong-expanded}
\end{equation}
since $\hbar\lambda=g/c$.  Its properties:
\begin{itemize}[leftmargin=*]
\item It is invariant under reparametrizations (all terms are of first order
  in $\dd/\dd\lambda$).
\item It is \emph{exactly} gauge invariant, with no total derivative:
  $D_\lambda\psi\to U D_\lambda\psi$ and $\psi^\dagger U^\dagger U\,D_\lambda\psi=\psi^\dagger D_\lambda\psi$.
\item In the abelian case $\psi=e^{\ii\theta}$, the internal term reduces to
  $-\hbar\dot\theta-\tfrac qcA_\mu\dot x^\mu$ with $-\hbar\dot\theta$ a total
  derivative, and one recovers \eqref{eq:S-EM}.
\item The internal term is first order in time derivatives (symplectic)
  with no internal energy; it is a Berry-phase-like term.
\end{itemize}
Varying $\psi^\dagger$ gives $\ii\hbar D_\lambda\psi=0$, i.e.\ \emph{parallel
transport of the colour spinor},
\begin{equation}
  \dot\psi=-\ii\lambda\,\dot x^\mu\mathcal A_\mu\,\psi ,
  \label{eq:psi-transport}
\end{equation}
which preserves $\psi^\dagger\psi$.  The colour charge then evolves as
\begin{equation}
  \dot Q^a=g\bigl(\dot\psi^\dagger T^a\psi+\psi^\dagger T^a\dot\psi\bigr)
  =\ii\lambda\,\dot x^\mu A^b_\mu\,g\,\psi^\dagger[T^b,T^a]\psi
  =\lambda f^{abc}\,\dot x^\mu A^b_\mu\,Q^c .
  \label{eq:Qdot}
\end{equation}
For $x^\mu$ one finds, from \eqref{eq:L-wong-expanded} with $\lambda=\tau$,
\begin{equation}
  \dot p_\mu=\frac1c\,Q^au^\nu\partial_\mu A^a_\nu,\qquad
  p_\mu=m u_\mu+\frac1c\,Q^aA^a_\mu .
\end{equation}
Expanding $\dot p_\mu$ and substituting \eqref{eq:Qdot},
\begin{align}
  m\dot u_\mu&=\frac1cQ^au^\nu\bigl(\partial_\mu A^a_\nu-\partial_\nu A^a_\mu\bigr)
  -\frac\lambda c\,f^{abc}A^a_\mu A^b_\nu u^\nu Q^c \nonumber\\
  &=\frac1c\,Q^aF^a_{\mu\nu}u^\nu ,
\end{align}
where the non-linear term of \eqref{eq:Fa} was recognized after a cyclic
relabelling of colour indices.  Collecting the two equations gives
\emph{Wong's equations}:
\begin{equation}
  \boxed{\ 
  m\,\frac{\dd u^\mu}{\dd\tau}=\frac1c\,Q^aF^{a\,\mu\nu}u_\nu,\qquad
  \frac{\dd Q^a}{\dd\tau}-\lambda f^{abc}\,u^\mu A^b_\mu\,Q^c=0\ }
  \label{eq:Wong}
\end{equation}
The second equation is $D_\tau Q^a=0$ with the adjoint covariant derivative
$(D_\mu Q)^a=\partial_\mu Q^a-\lambda f^{abc}A^b_\mu Q^c$.  Properties:
\begin{enumerate}[label=(\roman*),leftmargin=*]
\item \emph{Gauge covariance.}  $Q^aT^a\to U(Q^bT^b)U^{-1}$ (adjoint
  rotation), so $Q^aF^a_{\mu\nu}$ is gauge invariant and the world line is
  too.  Individual components $Q^a$ have no gauge-invariant meaning.
\item \emph{Casimirs.}  $D_\tau$ commutes with contraction by invariant
  tensors, so $Q^aQ^a$ and, for $SU(N)$, $d^{abc}Q^aQ^bQ^c$, etc.\ are
  constants of motion.  They label the orbit of $Q$ (Section~\ref{sec:orbit}).
\item \emph{Mass shell.}  $u_\mu\,\dd u^\mu/\dd\tau=0$ by antisymmetry of
  $F^a_{\mu\nu}$.  The energy changes only through the colour-electric field:
  $m\,\dd u^0/\dd\tau=\tfrac1c\,Q^aE^{a\,i}u^i$ with $E^{a\,i}=F^a_{0i}$.
\item \emph{Abelian limit.} One generator, $f=0$: $Q$ is constant and
  \eqref{eq:Wong} is \eqref{eq:Lorentz-cov} with $q=Q$.
\end{enumerate}

\subsection{Group variables and coadjoint orbits}\label{sec:orbit}

Write $\psi=g\,\psi_0$ with $g(\lambda)\in G$ (in the representation $R$)
and $\psi_0$ a fixed unit vector, $P_0=\psi_0\psi_0^\dagger$.  Then
\begin{equation}
  L_{\rm col}=\ii\hbar\,\tr\bigl(P_0\,g^{-1}D_\lambda g\bigr),\qquad
  D_\lambda g=\dot g+\ii\lambda\,\dot x^\mu\mathcal A_\mu g,\qquad
  Q^a=g\,\tr\bigl(T^a\,gP_0g^{-1}\bigr).
\end{equation}
The Lagrangian is unchanged (up to a total derivative) under
$g\to gh$ for $h$ in the stabilizer $H\subset G$ of the line through
$\psi_0$, so the physical internal space is the coset $G/H$: $\mathbb{CP}^{N-1}$
for the fundamental of $SU(N)$, $S^2$ for $SU(2)$.  In geometric language
$Q=\mathrm{Ad}^*_g\,Q_0$ moves on a \emph{coadjoint orbit}, and the first-order
term is the Kirillov--Kostant--Souriau symplectic potential on that orbit.
The Poisson bracket that follows from \eqref{eq:L-wong-expanded}, where the
momentum conjugate to $\psi$ is $\ii\hbar\psi^\dagger$ and hence
$\PB{\psi_i}{\psi^\dagger_j}=-\tfrac\ii\hbar\delta_{ij}$, is the
Lie--Poisson bracket
\begin{equation}
  \PB{Q^a}{Q^b}=\frac g\hbar\,f^{abc}Q^c=c\lambda\,f^{abc}Q^c .
  \label{eq:LiePoisson}
\end{equation}
Its Casimirs are central: they are the invariants fixed by the choice of
$Q_0$, i.e.\ by the representation.  A different but equivalent description
uses Grassmann colour variables $\theta_i$ with $Q^a=g\,\theta^\dagger T^a\theta$
(Balachandran et al., 1977; Brink et al., 1976), which leads to the same
Wong equations for the bilinears $Q^a$.

\subsection{Hamiltonian formalism}\label{sec:NA-ham}

The colour sector does not contribute to $\partial L/\partial\dot x^\mu$, so
\begin{equation}
  p_\mu=\frac{mc\,\dot x_\mu}{\sqrt{\dot x^2}}+\frac1c\,Q^aA^a_\mu,\qquad
  \pi_\mu\equiv p_\mu-\frac1c\,Q^aA^a_\mu,\qquad
  \pi_\mu\pi^\mu=m^2c^2 ,
\end{equation}
and the total Hamiltonian is again $H_T=e\,\Cc$ with
$\Cc=\tfrac12(m^2c^2-\pi_\mu\pi^\mu)$.  The phase space is
$(x^\mu,p_\mu)$ together with the colour charges, with brackets
\eqref{eq:PB-conv} and \eqref{eq:LiePoisson} and
$\PB{x^\mu}{Q^a}=\PB{p_\mu}{Q^a}=0$.  Then:
\begin{align}
  \dot x^\mu&=e\,\pi^\mu,\\
  \dot Q^a&=\PB{Q^a}{e\Cc}=-e\,\pi^\mu\PB{Q^a}{\pi_\mu}
  =\lambda f^{abc}\dot x^\mu A^b_\mu Q^c,
  \label{eq:Qdot-ham}\\
  \dot p_\mu&=\frac1c\,Q^a\dot x^\nu\partial_\mu A^a_\nu ,
\end{align}
where $\PB{Q^a}{\pi_\mu}=-\lambda f^{abc}A^b_\mu Q^c$.  These reproduce
\eqref{eq:Wong}.  The brackets of the gauge-covariant momenta are
\begin{equation}
  \boxed{\ \PB{\pi_\mu}{\pi_\nu}=-\frac1c\,Q^aF^a_{\mu\nu},\qquad
  \PB{\pi_\mu}{Q^a}=\lambda f^{abc}A^b_\mu Q^c\ }
  \label{eq:PB-NA}
\end{equation}
Indeed, 
$\PB{\pi_\mu}{\pi_\nu}=-\tfrac1cQ^a(\partial_\mu A^a_\nu-\partial_\nu A^a_\mu)
+\tfrac\lambda c\,f^{abc}A^a_\mu A^b_\nu Q^c$, and the last term is the
non-linear part of $-\tfrac1cQ^cF^c_{\mu\nu}$.  The Lorentz-force structure
of Section~\ref{sec:cov-ham} is therefore unchanged: \emph{the force is the
curvature of the kinetic momenta}, now weighted by the colour charge.
(The second bracket in \eqref{eq:PB-NA} is not gauge covariant by itself:
$Q^a$ and $\pi_\mu$ carry different $A$-dependent gauge transformations in
canonical variables.)

Solving $\Cc=0$ for $p^0$ with $A^a_0=\phi^a$, $\bA^a=(A^{a\,i})$ gives the
coordinate-time Hamiltonian
\begin{equation}
  H=Q^a\phi^a+c\sqrt{\Bigl(\bp-\frac1c\,Q^a\bA^a\Bigr)^2+m^2c^2},\qquad
  \dot f=\PB fH,
  \label{eq:H-NA}
\end{equation}
with the ordinary $\PB{x^i}{p^j}=\delta^{ij}$ and \eqref{eq:LiePoisson}.  The
corresponding Lagrangian is
\begin{equation}
  L=-\frac{mc^2}\gamma-Q^a\phi^a+\frac1c\,Q^a\,\bv\!\cdot\!\bA^a+L_{\rm col}.
\end{equation}
The colour charge enters exactly as $q$ does in
\eqref{eq:H-EM}, but it also \emph{evolves}, with
$\dd Q^a/\dd t=\lambda f^{abc}(c\,\phi^b-\bv\!\cdot\!\bA^b)Q^c$.

\subsection{SHP mechanics in a non-abelian field}

Combining \eqref{eq:L-SHP} and \eqref{eq:L-wong-expanded}, with the fifth
component acting like an extra ``time-like'' colour potential rotating $Q$:
\begin{equation}
  L=-\frac M2\,\dot x^2+\ii\hbar\,\psi^\dagger D_\tau\psi,\qquad
  D_\tau\psi=\dot\psi+\ii\lambda\bigl(\dot x^\mu\mathcal A_\mu+c\,\mathcal A_5\bigr)\psi,
\end{equation}
which gives
\begin{equation}
  M\ddot x_\mu=\frac1cQ^aF^a_{\mu\nu}\dot x^\nu+Q^aF^a_{\mu5},\qquad
  \dot Q^a=\lambda f^{abc}\bigl(\dot x^\mu A^b_\mu+c\,A^b_5\bigr)Q^c,
\end{equation}
\begin{equation}
  F^a_{\mu5}=\partial_\mu A^a_5-\frac1c\,\partial_\tau A^a_\mu-\lambda f^{abc}A^b_\mu A^c_5,\qquad
  K=-\frac1{2M}\pi_\mu\pi^\mu+Q^aA^a_5 .
\end{equation}
The fifth component both exerts a force on the event and rotates the colour
charge, and the ``mass'' $m^2c^2=M^2\dot x^2=2M(Q^aA^a_5-K)$ is exchanged
with the colour-fifth-potential energy, as in \eqref{eq:mass-SHP}.

\subsection{The geodesic (Kaluza--Klein) viewpoint}\label{sec:KK}

Both the Lorentz force and Wong's equations can be obtained as geodesic
motion in a higher-dimensional space.  In the abelian case consider
\begin{equation}
  L_5=-\frac M2\Bigl[\dot x^\mu\dot x_\mu-\bigl(\dot y+\kappa A_\mu\dot x^\mu\bigr)^2\Bigr],
\end{equation}
the geodesic Lagrangian of $\dd s_5^2=\eta_{\mu\nu}\dd x^\mu\dd x^\nu-(\dd y+\kappa A_\mu\dd x^\mu)^2$.
The coordinate $y$ is cyclic, so $p_y=M(\dot y+\kappa A_\mu\dot x^\mu)$ is
conserved.  Eliminating $y$ with the Routhian $R=L_5-p_y\dot y$,
\begin{equation}
  R=-\frac M2\,\dot x^2+\kappa\,p_y\,A_\mu\dot x^\mu-\frac{p_y^2}{2M},
\end{equation}
which is \eqref{eq:LS-free} with minimal coupling, provided
$\kappa p_y=-q/c$: the charge is the conserved momentum of the extra
dimension.  This is the origin of the observation (Kaluza, Klein) that
electric charge is momentum in a compact fifth dimension.  For a
non-abelian group, replace the circle by $G$: with the Maurer--Cartan forms
$\sigma^a$ on $G$ and a bi-invariant metric $\gamma_{ab}$, the metric
\begin{equation}
  \dd s^2=\eta_{\mu\nu}\dd x^\mu\dd x^\nu-\gamma_{ab}\bigl(\sigma^a+\kappa A^a_\mu\dd x^\mu\bigr)\bigl(\sigma^b+\kappa A^b_\nu\dd x^\nu\bigr)
\end{equation}
on $M\times G$ has geodesics that project to solutions of Wong's equations,
and the conserved momenta of the $G$-action are the colour charges
(Kerner, 1968; see also Sternberg 1977 and Weinstein 1978 for the
symplectic version of this statement).  We state this without proof;
the normalization of $\kappa$ and the sign conventions must be matched to
\eqref{eq:D-F}.  If the extra direction is taken \emph{null} instead of spacelike, the same
Routh reduction gives Galilean (Bargmann) dynamics, with mass playing the
role of $p_y$ (null reduction).

\subsection{Two exactly solvable examples}

\paragraph{Pure gauge.}  If $\mathcal A_\mu=\tfrac\ii\lambda(\partial_\mu U)U^{-1}$
then $\mathcal F_{\mu\nu}=0$, there is no force, and $D_\mu(U\psi_0)=0$ shows
that $\psi(\tau)=U(x(\tau))\psi_0$: the colour is merely rotated by the
adjoint action of the gauge function, and the world line is a free
straight line.

\paragraph{Abelian-dominant field in $SU(2)$.}  Let $A^a_\mu=\delta^{a3}a_\mu(x)$.
Then $F^a_{\mu\nu}=\delta^{a3}f_{\mu\nu}$ with $f_{\mu\nu}=\partial_\mu a_\nu-\partial_\nu a_\mu$
(the non-linear term vanishes since $f^{a33}=0$).  Equations \eqref{eq:Wong} give
$\dot Q^3=0$ and
\begin{equation}
  \frac{\dd}{\dd\tau}(Q^1+\ii Q^2)=\ii\lambda\,u^\mu a_\mu\,(Q^1+\ii Q^2)
  \ \Rightarrow\ 
  Q^1+\ii Q^2=\bigl(Q^1+\ii Q^2\bigr)_0\exp\Bigl[\ii\lambda\!\int a_\mu\dd x^\mu\Bigr].
\end{equation}
The particle moves exactly like an electromagnetic particle of charge
$Q^3$ in the field $f_{\mu\nu}$---all results of Sections~\ref{sec:em}--\ref{sec:examples}
apply with $q\to Q^3$---while the transverse colour components precess by
the Wilson-line phase.  Genuinely non-abelian effects require non-commuting
components of $\mathcal A_\mu$ (Problem~\ref{prob:const}).

\subsection{Remark on the quantum theory}

The same world line data arises from the multi-component Klein--Gordon equation
$\bigl[D_\mu D^\mu+m^2c^2/\hbar^2\bigr]\psi=0$ with $D_\mu$ of
\eqref{eq:D-F} acting on $\psi(x)\in\mathbb C^d$, and Wong's equations are
its classical limit when the colour degree of freedom is treated as an
independent classical variable (large representation, coherent-state
limit).  The derivation is more delicate than in the abelian case, since the
eikonal equation is matrix valued; see the notes.

% =====================================================================
\section{Summary and comparison}\label{sec:summary}
% =====================================================================

\begin{itemize}[leftmargin=*]
\item The \emph{coordinate-time} Lagrangian and Hamiltonian
  \eqref{eq:L-EM-t}, \eqref{eq:H-EM}, \eqref{eq:H-NA} are regular and best
  for the nonrelativistic limit, perturbation theory, and numerical
  integration, where one wants an explicit $H(\bx,\bp,t)$ and a symplectic
  integrator.

\item The \emph{covariant} (square-root or einbein) forms exhibit
  reparametrization invariance as a gauge symmetry generated by the
  first-class constraint $\Cc\approx0$.  They are the natural starting
  point for quantization (Klein--Gordon from $\Cc$), for the BRST
  formalism, and for generalizations to strings and spinning particles.

\item The \emph{SHP} form removes the constraint at the cost of an extra
  parameter and an extra potential $A_5$; the mass becomes dynamical.

\item The \emph{curvature of kinetic momenta}, $\PB{\pi_\mu}{\pi_\nu}=-\tfrac1cQ^aF^a_{\mu\nu}$,
  is the same statement in the abelian and non-abelian cases; the new
  physics in the non-abelian case is that $Q^a$ evolves on a coadjoint
  orbit under the Lie--Poisson bracket, and only the Casimirs of $Q$ are
  gauge-invariant.

\item The \emph{geodesic} picture explains \emph{why} charge is conserved
  (a cyclic coordinate) and why colour precesses (the G-action), and
  connects the particle dynamics to Kaluza--Klein theory.
\end{itemize}

% =====================================================================
\section*{Problems}
% =====================================================================

\begin{problem}
Show that $S_e$ in \eqref{eq:Se-free} is invariant under $\lambda\to\lambda'(\lambda)$
if $e'(\lambda')=e(\lambda)\,\dd\lambda/\dd\lambda'$, and that the
Hamiltonian derived from $S_e$ is $e\Cc$.
\end{problem}

\begin{problem}
Solve the Hamilton--Jacobi equation \eqref{eq:HJ-EM} for a uniform electric
field by separation of variables and recover \eqref{eq:hyperbolic}.
\emph{Hint:} take $\phi=-Ex$, $\bA=0$, $S=-Et+W(x)$.
\end{problem}

\begin{problem}
For crossed fields $\bE\perp\bB$ with $|\bE|<|\bB|$, show that the drift
velocity is $\bv_d=c\,\bE\times\bB/B^2$.  \emph{Hint:} boost to a frame with
$\bE'=0$.
\end{problem}

\begin{problem}
Verify \eqref{eq:Volkov-classical} by substitution and show that a particle
initially at rest in a linearly polarized plane wave acquires the
Lorentz factor $\gamma=1+\tfrac{q^2}{2m^2c^4}|\mathbf A-\mathbf A_0|^2$.
Find its mean drift velocity.
\end{problem}

\begin{problem}
Prove the bracket \eqref{eq:PB-NA} and use it to rederive Wong's equations
from $H_T=e\Cc$.
\end{problem}

\begin{problem}
Show that for $SU(3)$ the cubic invariant $d^{abc}Q^aQ^bQ^c$ is a constant of
motion of Wong's equations.  \emph{Hint:} the invariance identity for the
symmetric tensor $d^{abc}$ is $f^{ade}d^{ebc}+f^{bde}d^{aec}+f^{cde}d^{abe}=0$.
\end{problem}

\begin{problem}
Starting from \eqref{eq:EOM-SHP} derive $\dd(\tfrac12M\dot x^2)/\dd\tau=q\,\dot x^\mu F_{\mu5}$,
and show that $K$ is conserved if and only if $\partial_\tau$ of the
fields vanishes.
\end{problem}

\begin{problem}
Check that Wong's equations are covariant under finite gauge
transformations $U(x)$: if $(x,Q^a,A^a_\mu)$ solves them, so does
$(x,Q'^a,A'^a_\mu)$ with $Q'T=UQTU^{-1}$ and $A'$ given by \eqref{eq:gauge-NA}.
\end{problem}

\begin{problem}\label{prob:const}
For $SU(2)$ take constant potentials $A^1_x=\alpha$, $A^2_y=\beta$ and all
others zero.  Show that $F^3_{xy}=-\lambda\alpha\beta$ is a constant
colour-magnetic field with non-commuting potentials, and analyse the Wong
motion for a heavy particle.  Why is the particle \emph{not} simply a charge
$Q^3$ in a uniform magnetic field?
\end{problem}

% =====================================================================
\section*{Notes and further reading}
% =====================================================================

\paragraph{Classical textbook treatments.}  The covariant electromagnetic
case is treated in L.~D.\ Landau and E.~M.\ Lifshitz, \emph{The Classical
Theory of Fields}, and J.~D.\ Jackson, \emph{Classical Electrodynamics}
(note that Jackson uses Gaussian units in the older editions and the
$(+,-,-,-)$ metric).

\paragraph{Constraints, einbein.}  P.~A.\ M.\ Dirac, \emph{Lectures on
Quantum Mechanics} (1964); M.\ Henneaux and C.\ Teitelboim,
\emph{Quantization of Gauge Systems} (1992), for the constraint analysis
of reparametrization-invariant systems.  A.~M.\ Polyakov, \emph{Gauge
Fields and Strings} (1987), for the einbein form.  L.\ Brink, S.\
Deser, B.\ Zumino, P.\ Di~Vecchia and P.\ Howe, Phys.\ Lett.\ B64 (1976)
435, for the spinning particle.

\paragraph{SHP.}  E.~C.~G.\ Stueckelberg, Helv.\ Phys.\ Acta 14 (1941) 322, 588;
L.~P.\ Horwitz and C.\ Piron, Helv.\ Phys.\ Acta 46 (1973) 316;
L.~P.\ Horwitz, \emph{Relativistic Quantum Mechanics} (Springer, 2015).
The literature on SHP uses the metric $(-,+,+,+)$, so $K=+p^2/2M$ there.

\paragraph{Non-abelian particles.}  S.~K.\ Wong, Nuovo Cimento A65 (1970)
689; R.\ Kerner, Ann.\ Inst.\ H.\ Poincar\'e 9 (1968) 143;
A.~P.\ Balachandran, S.\ Borchardt and A.\ Stern, Phys.\ Rev.\ D17 (1978) 3247;
A.~P.\ Balachandran, P.\ Salomonson, B.-S.\ Skagerstam and J.-O.\ Winnberg,
Phys.\ Rev.\ D15 (1977) 2308 (Grassmann colour variables);
S.\ Sternberg, Proc.\ Natl.\ Acad.\ Sci.\ USA 74 (1977) 5253;
A.\ Weinstein, Lett.\ Math.\ Phys.\ 2 (1978) 417;
R.\ Montgomery, Lett.\ Math.\ Phys.\ 8 (1984) 59;
B.\ P.\ Kosyakov, \emph{Introduction to the Classical Theory of Particles
and Fields} (Springer, 2007).  For the classical limit of the quantum
colour dynamics see U.\ Heinz, Phys.\ Rev.\ Lett.\ 51 (1983) 351.

\paragraph{Feynman's derivation.}  F.~J.\ Dyson, Am.\ J.\ Phys.\ 58 (1990) 209.

\medskip
\noindent\emph{The bibliographic details above are given from memory and
should be checked against the primary sources before publication.}

\end{document}
